All three are feed-forward neural networks: information passes once from input to output, with no
internal memory carried between time steps. They differ from one another in how the layers are
arranged, described below.

\subsection*{N-BEATS}

A deep stack of fully connected layers with backward and forward residual links. It builds the
forecast as a sum of learned basis functions, which is what makes it partly interpretable. Designed
for univariate point forecasting and benchmarked on the M3, M4 and TOURISM competition datasets.

\noindent Implemented in Darts as \href{https://unit8co.github.io/darts/generated_api/darts.models.forecasting.nbeats.html}{\texttt{NBEATSModel}};
see \citet{oreshkin2020nbeats}.


\subsection*{TSMixer}

Dense layers only, applied alternately along the time axis and across the predictor variables, so
the model mixes temporal structure and cross-variable structure in turn. Built for the multivariate
case.

\noindent Implemented in Darts as \href{https://unit8co.github.io/darts/generated_api/darts.models.forecasting.tsmixer_model.html}{\texttt{TSMixerModel}};
see \citet{chen2023tsmixer}.


\subsection*{TiDE}

An encoder--decoder, again dense layers only. It compresses the history and the predictors into a
single code, then decodes the whole 36-month path from it in one pass.

\noindent Implemented in Darts as \href{https://unit8co.github.io/darts/generated_api/darts.models.forecasting.tide_model.html}{\texttt{TiDEModel}};
see \citet{das2023tide}.

